\section{Recommendations}\label{recommendations}

The following recommendations follow directly from the evidence limits
identified in the drainage, software-dependency, and electrical-network
cases:

\begin{enumerate}[label=\arabic*.,leftmargin=*]
  \item Obtain authorized electrical measurements covering the four
  requested November 13 event groups and execute the preregistered
  event-window analysis.
  \item Evaluate domain-specific rankings against independent outcomes
  that were not used to construct the model inputs.
  \item Obtain provenance-complete operational drainage observations so
  that the controlled and SWMM-based findings can be compared with field
  conditions rather than being interpreted as field validation.
  \item Investigate ancestry-aware or dependence-aware aggregation for
  reconvergent paths while retaining the current v1 formulation as the
  comparison baseline.
  \item Calibrate or justify heterogeneous transmission and
  susceptibility functions separately within each domain.
  \item Develop explicit extensions for cyclic, bidirectional, and
  temporal networks rather than silently applying the DAG evaluator
  outside its assumptions.
  \item Repeat the cross-domain protocol on additional datasets to test
  out-of-sample applicability without overstating generalizability.
\end{enumerate}

These recommendations preserve the distinction between the model's
validated computational behavior and the empirical claims that still
require domain-specific evidence. GBBRPM v1 is not a universal predictor,
a hydraulic simulator, or a completed electrical event model; subsequent
work should retain those interpretation boundaries while strengthening
field evidence and dependence handling.
